\clearpage
\fancyhead[L]{\small VOYAGE DE DENTRECASTEAUX.}
\hypertarget{native-519-lower}{}\label{native-519-lower}
\begin{center}\small Göttingische gelehrte Anzeigen. Stück 167. Seite 1657..1663. 1809 October 21.\\[3pt]\rule{.72\linewidth}{.3pt}\end{center}

\emph{Voyage de} \textsc{Dentrecasteaux}\emph{, envoyé à la recherche de la Pérouse. Publié par ordre de Sa Majesté l'Empereur et Roi.} \emph{Rédigé par} M. \textsc{de Rossel}\emph{, ancien capitaine de vaisseau.} \emph{Tome second.} \emph{A Paris} 1808. 692 Seiten in Quart. Paris. Der erste Theil dieses Werks, welcher die Reise selbst enthält, ist bereits (oben S. 1065) in diesen Blättern angezeigt. Dieser zweite Band ist ganz den astronomisch-nautischen Beobachtungen gewidmet und zerfällt in zwei Hauptabtheilungen, wovon die erstere die Beschreibung der Instrumente und allgemeine Untersuchungen über den damit zu erreichenden Grad von Genauigkeit und die Umstände, wo die Beobachtungen am vortheilhaftesten anzustellen sind, enthält; die andere hingegen die während der Reise auf der Fregatte La Recherche, zum Theil auch die auf der andern Fregatte, L'Esperance, gemachten astronomischen Beobachtungen im grössten Detail, tabellarisch geordnet, aufstellt. Da die für die Sicherheit der Schifffahrt höchst wichtigen astronomischen Beobachtungen und das Geschäft, Breiten- und Längenbestimmungen daraus abzuleiten, meistens Personen zufällt, die keine eigentlich mathematische Bildung erhalten, sondern sich gewöhnlich wenig mehr als eine mechanische Fertigkeit in Beobachtungen und Rechnungen erworben haben, so ist der Tractat, welcher die erstere Abtheilung ausmacht. von einem solchen sachverständigen und erfahrnen Beobachter, wie der Capitän \textsc{Rossel}, gewiss eine sehr schätzbare Arbeit, und das Studium desselben angehenden Seemännern, zur Schärfung der Beurtheilung des 'den Beobachtungen beizulegenden Werths, sehr zu empfehlen. Auch Liebhaber, die zu Lande beobachten, werden denselben mit Nutzen lesen. Nur wäre zu wünschen, dass dieser Tractat nicht einen Theil eines so theuern Werks ausmachte, mit dem er eigentlich in keiner unmittelbaren Verbindung steht, oder dass er doch wenigstens durch einen besondern Abdruck mehreren Lesern zugänglich gemacht würde. Auch hätten wir wohl gewünscht, dass durch kleinere Unterabtheilungen für eine bequemere Uebersicht gesorgt wäre: es sind zwar Paragraphen da, aber Paragraphen von 119 S. Länge. Von den zwei Capiteln, in welche die Abhandlung zerfällt, handelt das erste von den Instrumenten, womit die beiden Fregatten ausgerüstet waren, und erläutert zugleich die jedem Instrumente eigenthümlichen Vorzüge. Jede Fregatte hatte, ausser mehreren \textsc{Bordai}'schen Reflexionskreisen, einen astronomischen Repetitionskreis für den Gebrauch zu Lande, eine Seeuhr, eine astronomische Uhr, einen Zähler, ein achromatisches Fernrohr, eine Inclinationsbous\SourceLeaf{526}sole nach einer neuen Einrichtung, und einen Azimuthalcompass. Mit den Inclinations-Boussolen wurden auch an mehreren Orten Versuche über die Dauer der Oscillationen der Nadel im magnetischen Meridiane angestellt, und \textsc{Rossel} fügt eine Tafel zur Reduction der Oscillationen auf unendlich kleine bei; der ganze Apparat befand sich in einem mit Glasscheiben verschlossenen Gehäuse, welches aber bei diesen Versuchen geöffnet werden musste, daher \textsc{Rossel} glaubt, dass zuweilen die Bewegung der Luft die Versuche etwas ungewiss gemacht haben könne. Endlich wird noch die Einrichtung eines astronomischen Zeltes beschrieben, dergleichen jedes Schiff eins führte, um bei den Beobachtungen zu Lande gebraucht zu werden. — Das zweite Kapitel entwickelt die Natur und Grenzen der Fehler, die von den Werkzeugen, von den Beobachtungen selbst oder von den Tafeln herrühren, und die Mittel, sie zu verbessern oder wenigstens nach Möglichkeit zu vermindern. Die verschiedenen Gattungen von Beobachtungen werden hier in 9 Paragraphen abgehandelt. §. I. über die Zeitbestimmung durch einzelne Höhen. Der Verf. macht hier unter andern auf den Einfluss der bei Bestimmung der Depression des Horizonts begangenen Fehler aufmerksam und empfiehlt daher, immer in einer so grossen Höhe über dem Spiegel der Meeresfläche, als es thunlich ist, zu beobachten: offenbar gilt dieselbe Bemerkung bei allen zur See beobachteten Höhenwinkeln. §. II. Breitenbestimmungen durch Meridianhöhen. Sonnenhöhen nahe beim Zenith hält \textsc{Rossel} nur auf 2..3 Minuten genau: der Grund, warum er unter solchen Umständen die Messung für weniger zuverlässig hält, weil nemlich beim Wiegen des Reflexionskreises das Bild der Sonne sich nur sehr langsam vom Seehorizonte trennt, und daher der eigentliche Verticalkreis der Berührung nicht genau bestimmt werden kann, ist uns nicht recht einleuchtend. Sternhöhen zur See zu nehmen, wollte nie gelingen, auch bei den günstigsten Umständen waren Fehler bis zu 5 Minuten nicht zu vermeiden: vielleicht liegt das zum Theil an der Unvollkommenheit des optischen Theils der französischen Werkzeuge, wo die Fernröhre eine gar zu kleine Oeffnung haben. Breitenbestimmungen durch Mondshöhen setzt \textsc{Rossel} bei Tage auf 3, bei Nacht auf 4 Minuten genau: dabei ist aber 1 Minute als Fehler der Declination des Mondes aus den Ephemeriden eingerechnet, welcher sich leicht vermeiden liesse, wenn man sich die Mühe nicht verdriessen lässt, die Declination aus der Länge und Breite zu berechnen, welche in den Ephemeriden auf Secunden angegeben zu werden pflegen. §. III. Breitenbestimmung aus zwei Sonnenhöhen ausser dem Meridian. Mit Recht werden die indirecten Auflösungsmethoden vorzugsweise empfohlen, wenn auch die directe Auflösung nicht nothwendig so sehr weitläuftige und beschwerliche Rechnungen und die verwickelte Unterscheidung vieler Fälle erfordert, wie \textsc{Rossel} glaubt. Ueber die Auflösung selbst findet man zwar hier nichts Neues; aber sehr umständlich ist der Einfluss der Fehler in den Beobachtungen, in der Schätzung der Fortbewegung des Schiffes zwischen den Beobachtungen, und im Gange der Uhr erörtert. Das Endresultat ist, dass man durch dies Verfahren auch im schlimmsten Falle, wenn alle Fehler in Einem Sinne conspiriren, doch immer die Breite auf 4 bis 5 Minuten genau bestimmen könne, wenn man bei der Auswahl der Beobachtungen die nöthigen Vorsichtsregeln anwendet. Was über die grössere oder geringere Genauigkeit der verschiedenen indirecten Rechnungsmethoden gesagt wird, dürfte bei den in der Ausübung vorkommenden Fällen ziemlich überflüssig sein. §. IV. Beobachtung des Azimuths der Sonne zur Bestimmung der Abweichung der Magnetnadel. \textsc{Rossel} macht hier aufmerksam darauf, dass zuweilen die Abweichung an der Küste am Bord des Schiffs von der in geringer Entfernung davon am Lande gefundenen beträchtlich verschieden ausfällt, wovon die Beobachtungen bei Teneriffa, welche wir weiter unten anführen werden, ein merkwürdiges Beispiel geben. §. V. Längenbestimmungen aus gemessenen Abständen des Mondes von der Sonne und von Fixsternen. Man findet hier manche schätzbare praktische Bemerkung, die als die Frucht einer langen Erfahrung um so mehr Gewicht hat. Denn gewiss mit Recht sagt der Verf.: Tout ce qui peut faciliter une observation aussi délicate, est précieux; nous n'en con\SourceLeaf{527}noissons pas qui exige plus de dispositions naturelles: aussi nous sommes bien convaincus que, lorsque les savans auront fait, pour la perfection de cette méthode, tout ce que leur génie pourra leur inspirer, le talent de l'observateur aura encore beaucoup à y ajouter. — Bei Berechnung der Mondsparallaxe nimmt \textsc{Rossel}, zum Behuf der Reduction der gemessenen Abstände, auf die sphäroidische Gestalt der Erde nur in so fern Rücksicht, als die Horizontalparallaxe nach der Breite des Beobachtungsortes kleiner ist als die Aequatorealparallaxe, und vernachlässigt die Neigung der Verticale gegen den zum Mittelpunkte der Erde gehenden Halbmesser. Dies hat nun freilich bei Beobachtungen, wie sie mit gewöhnlichen französischen Spiegelkreisen gemacht werden können, nicht viel zu bedeuten: allein den Einfluss dieses letztern Umstandes hat \textsc{Rossel}, nach unserer Meinung, nicht richtig beurtheilt. Er nimmt nemlich nur darauf Rücksicht, dass die Parallaxe nach einer falschen \emph{Höhe} berechnet wird, und vergisst, dass auch die \emph{Richtung} der Parallaxe dadurch afficirt wird. Daher halten wir diesen Schluss, das Resultat werde am meisten durch diese nicht ganz richtige Rechnungsart entstellt, wenn der Mond nahe bei der Culmination sei, für irrig, indem in der Regel der Fehler gerade am grössten wird, wenn der Mond sich nahe beim ersten Vertical befindet. §. VI. Längenbestimmungen durch Seeuhren. §. VII. Ueber die Art, wie die durch Mondsdistanzen bestimmten Längen mit den durch Seeuhren erhaltenen Längenunterschieden verbunden werden müssen. §. VIII. Längenbestimmungen aus Sternbedeckungen und Finsternissen der Jupiterstrabanten. §. IX. Bestimmungen der Azimuthe terrestrischer Gegenstände (relevemens).

Die zweite Abtheilung des Werks enthält das Beobachtungs-Journal selbst, welchem eine Tafel der Fehler der \textsc{Mason}'schen Mondstafeln nach gleichzeitigen Greenwicher Beobachtungen, und die daraus abgeleiteten Verbesserungen der während der Reise gemachten Längenbestimmungen vorausgeschickt wird. Von den für die Geographie des Südmeers sehr wichtigen Resultaten der Beobachtungen hier Etwas auszuheben, würde gegen den Plan dieser Gel. Anz. sein, auch um so unnöthiger, da man einen sehr vollständigen Auszug der geographischen Bestimmungen bereits in der Monatl. Correspondenz findet. Inzwischen wird es doch vielleicht manchem Leser angenehm sein, wenn wir hier die wenigen an der Magnetnadel gemachten Bestimmungen mittheilen, welche in jener Zeitschrift übergangen sind. Zu \emph{Brest} fand man am 20. September 1791 die Neigung der Nadel gegen Norden $71^\circ\,30'$, die Dauer einer unendlich kleinen Schwingung $2^{\mathrm s}02$. Zu \emph{Santa Cruz} auf Teneriffa war am 21. October 1791 die Neigung gegen Norden $62^\circ\,25'$; die mittlere Schwingungsdauer $2^{\mathrm s}081$; die Declination der Nadel am Bord der Recherche $18^\circ\,7'$ N. W., beim Observatorium $21^\circ\,33'$, und auf dem Molo $23^\circ\,43'$. Im \emph{Port du Nord} auf Van Diemensland den 11. Mai 1792 die Inclination $70^\circ\,50'$ südlich, die mittlere Schwingungsdauer $1^{\mathrm s}869$, die Abweichung $8^\circ\,1'\,12''$ N. O. Auf der Insel \emph{Amboina} den 9. October 1792 Neigung $20^\circ\,37'$ südlich, Schwingungsdauer $2^{\mathrm s}403$. Im \emph{Port au Sud} auf Van Diemensland den 7. Februar 1793 die Abweichung mit der einen Nadel $2^\circ\,33'\,47''$ N. O., mit der andern $3^\circ\,10'\,24''$; die Neigung $72^\circ\,22'$ und $70^\circ\,48'$ südlich; die Schwingungsdauer $1^{\mathrm s}8498$ und $1^{\mathrm s}817$. In \emph{Sourabaya} auf der Insel Java den 9. Mai 1794 die Neigung $25^\circ\,40'$ südlich, und die Schwingungsdauer mit der einen Nadel $2^{\mathrm s}429$, mit der andern $2^{\mathrm s}485$. Aus den verschiedenen Beobachtungen der Dauer der Oscillationen scheint also hervorzugehen, dass die Intensität der magnetischen Kraft unter dem Aequator beträchtlich schwächer ist, als in höheren Breiten.
